\section{Where the phase potential is monotone}\label{sec:monotone}

By Corollary~\ref{cor:increase} the phase potential is not monotone on all
complexes. This section treats two families: complexes with cohomology in
degree zero whose acyclic part is at one vertex of an arbitrary quiver,
where the monotonicity known for $K_m$ \cite[Theorem~7.6]{wue26b} persists,
and complexes over $K_m$ with acyclic parts at both vertices, where it
persists for ranks at most one and fails in general for larger ranks
(Theorem~\ref{thm:twosidedfailure}, Corollary~\ref{cor:classification}). We
use $\cD$ and $\cD_i^n$ from Section~\ref{sec:rate}, $K_\Delta$, $\mathcal J$
and the divided differences $\mathsf f[\,\cdot\,]$ from
Section~\ref{sec:setting}, $J_L(M)$, $M_\pm$ and the inequality
\eqref{eq:kernel} from Theorem~\ref{thm:imported}, and the terms of
Definition~\ref{def:phase}.

\subsection{An acyclic part at one vertex}
Let $Q$ be a finite quiver, $i_0$ a vertex without loops, and consider
complexes $\mathsf E$ with $E_i=E_i^0$ for $i\neq i_0$ and
$E_{i_0}=V\oplus U$ in degrees $0,1$, with $d_{i_0}=D:V\to U$ surjective and
$r=\dim U\geq1$. The chain-map condition says that every arrow into $i_0$
maps into $\ker D$, and is empty for the arrows out of $i_0$. The cohomology
is a representation $\mathcal H$ of $Q$ in degree $0$, with $\ker D$ at $i_0$;
we assume $\mathcal H\neq0$. Then
$Z=Z(\mathsf E)=\sum_iz_i\dim\mathcal H_i=|Z|e^{\ii\varphi}$ with
$\varphi\in(0,\pi)$. Let $\sigma=\sin\varphi$, $\alpha_i=\arg z_i$,
$c_0=|z_{i_0}|\sin(\alpha_{i_0}-\varphi)$, $B=D^\dagger D$, $L=DD^\dagger$,
\[
 C_{\mathrm{out}}=\sum_{a:i_0\to j}\rho(a)^\dagger\rho(a),\qquad
 G_{\mathrm{in}}=\sum_{a:k\to i_0}\rho(a)\rho(a)^\dagger,\qquad
 J=J_L(DC_{\mathrm{out}}D^\dagger),
\]
and $R(b)=(1-e^{-b})/b$.

\begin{theorem}[Acyclic part at one vertex]\label{thm:onevertex}
For the complexes just described:
\begin{enumerate}
\item[(a)] The blocks of $\Omega$ are
 \[
 \Omega_V=z_{i_0}e^{-B}+K_B(C_{\mathrm{out}})-G_{\mathrm{in}},\qquad
 \Omega_U=J-z_{i_0}e^{-L},
 \]
 and for $j\neq i_0$
 \[
 \Omega_j=z_j+\sum_{a:j\to\cdot}\rho(a)^\dagger\rho(a)
 -\sum_{\substack{a:k\to j\\k\neq i_0}}\rho(a)\rho(a)^\dagger
 -\sum_{a:i_0\to j}\rho(a)R(B)\rho(a)^\dagger .
 \]
\item[(b)] If $\alpha_{i_0}\geq\varphi$, then for every metric
 \[
 \cD\geq\varphi c_0\tr e^{-L}+\sigma\bigl(\varphi\tr J-\pi\tr J_-\bigr)
 \geq\varphi c_0\tr e^{-L}+\sigma\bigl(\varphi-\tfrac{2\pi}5\bigr)\tr J .
 \]
 Hence the phase potential is nonincreasing along the flow if
 $\alpha_{i_0}\geq\varphi$ and either $\varphi\geq\frac{2\pi}5$ or $r=1$, and
 strictly decreasing if moreover $\alpha_{i_0}>\varphi$.
\item[(c)] Let instead $E_i=E_i^0$ for $i\neq i_0$ and
 $E_{i_0}=U\oplus V$ in degrees $-1,0$, with $d_{i_0}=D:U\to V$ injective,
 $r=\dim U\geq1$, and nonzero cohomology, so that
 $\varphi=\arg Z(\mathsf E)\in(0,\pi)$. If $\alpha_{i_0}\leq\varphi$ and
 either $\varphi\leq\frac{3\pi}5$ or $r=1$, the phase potential is
 nonincreasing along the phase flow.
\end{enumerate}
\end{theorem}

\begin{proof}
(a) Proposition~\ref{prop:coefficient}. For $a:i_0\to j$,
$M_a=-\rho(a)D^\dagger:U\to E_j$. For $a:k\to i_0$, $M_a=d_{i_0}^\dagger\rho(a)=0$,
because $\rho(a)$ takes values in $V$, on which $d_{i_0}^\dagger$ vanishes.
For arrows not meeting $i_0$, $M_a=0$. So at $i_0$, $\Delta=B\oplus L$ and
$\Gamma=(z_{i_0}+C_{\mathrm{out}}-G_{\mathrm{in}})\oplus z_{i_0}$;
$K_B(G_{\mathrm{in}})=G_{\mathrm{in}}$ because $G_{\mathrm{in}}$ is supported on
$\ker B$; and $M_a^\dagger M_a=D\rho(a)^\dagger\rho(a)D^\dagger$ is supported on
$U$, with $\Delta_j=0$. At $j\neq i_0$, $\Delta_j=0$ and the mixed term is
$\rho(a)D^\dagger\bigl[\int_0^1(1-x)e^{-xL}dx\bigr]D\rho(a)^\dagger
=\rho(a)(1-R(B))\rho(a)^\dagger$.

(b) All blocks other than $U$ are in degree $0$, with phases in $(0,\pi)$,
and contribute nonnegative terms to $\cD$. Let $(u_k)$ be an orthonormal
eigenbasis of $\Phi_U$ with eigenvalues $\psi_k\in(-\pi,0)$,
$q_k=\inner{u_k}{|\Omega_U|u_k}$, $\iota_k=\inner{u_k}{Ju_k}$ and
$e_k=\inner{u_k}{e^{-L}u_k}>0$. Since
$\inner{u_k}{\Omega_Uu_k}=q_ke^{\ii\psi_k}$,
$q_k\sin(\psi_k-\varphi)=\Ima(e^{-\ii\varphi}\inner{u_k}{\Omega_Uu_k})
=-(\sigma\iota_k+c_0e_k)$. If $\sigma\iota_k+c_0e_k\geq0$, then
$\psi_k-\varphi\in[-\pi,-\varphi)$ and the summand
$q_k(\psi_k-\varphi)\sin(\psi_k-\varphi)$ of the contribution of $U$ is at
least $\varphi(\sigma\iota_k+c_0e_k)$; otherwise
$\psi_k-\varphi\in(-\pi-\varphi,-\pi)$ and the summand is at least
$(\pi+\varphi)(\sigma\iota_k+c_0e_k)$. In both cases it is at least
$\varphi(\sigma\iota_k+c_0e_k)-\pi\sigma(\iota_k)_-$, because $c_0e_k\geq0$;
here $(\iota_k)_-=\max(-\iota_k,0)$. Summing, with
$\sum_k(\iota_k)_-\leq\tr J_-$, gives the first inequality; the second is
\eqref{eq:kernel}, since $DC_{\mathrm{out}}D^\dagger\geq0$. For $r=1$,
$J\geq0$. If $\alpha_{i_0}>\varphi$, then $c_0>0$ and $\cD>0$. (c) follows from
(b) by Theorem~\ref{thm:duality}: the dual complex has an acyclic part at
$i_0$ in degrees $0,1$ with surjective differential, the charge
$-\overline{z_{i_0}}$ of argument $\pi-\alpha_{i_0}$ at $i_0$, and central
charge of argument $\pi-\varphi$ (Remark~\ref{rem:injective}).
\end{proof}

For $Q=K_m$, $i_0=0$ and cohomology at both vertices, $\varphi$ lies between
$\alpha_1$ and $\alpha_0$, and the case $\alpha_0>\alpha_1$ of (b) is
\cite[Proposition~7.4 and Theorem~7.6]{wue26b}; here equality and cohomology
at one vertex are
allowed. For $r\geq2$ a lower bound on $\varphi$ cannot be dispensed with:
for $K_2$, the charges $z_0=10^{-6}(-1+\ii)$, $z_1=1+\frac\ii{500}$ and
cohomology of dimension vector $(1,1)$, where $\varphi=0.0020\ldots$, there
is a complex with $r=2$ and a metric with $\cD<0$
\cite[Proposition~10.2]{wue26b}. The
hypothesis on $\alpha_{i_0}$ cannot be dropped: by
Proposition~\ref{prop:contractible} a contractible summand at $i_0$ in
degrees $0,1$ contributes
$\pi|z_{i_0}|\tr e^{-ff^\dagger}\sin(\alpha_{i_0}-\varphi)$ to $\cD$, and
Corollary~\ref{cor:increase} is an instance with a negative sign. Decay of
the acyclic part and convergence on the cohomology
(Definition~\ref{def:phase}(v)) for the complexes of
Theorem~\ref{thm:onevertex} are the subject of \cite{wue26b}, for $K_m$ with
a surjective differential at the vertex $0$, and of \cite{wue26c}, for
quivers without oriented cycles.

\subsection{Acyclic parts at both vertices of a generalized Kronecker quiver}
Let $Q=K_m$. Consider the three-term complexes
\begin{equation}\label{eq:twosided}
 (0\to U_1)\ \xrightarrow{\ D_1\ }\ (V\xrightarrow{A_j}W)\ \xrightarrow{\ D_0\ }\ (U_0\to0)
 \qquad\text{in degrees }-1,0,1,
\end{equation}
with $D_1:U_1\to W$ injective and $D_0:V\to U_0$ surjective; the chain-map
conditions are empty. The cohomology is the representation
$\ker D_0\to W/D_1U_1$ in degree $0$. Put $r_0=\dim U_0$, $r_1=\dim U_1$,
$n_0=\dim\ker D_0$, $n_1=\dim W-r_1$, assume $n_0+n_1\geq1$, and let
$Z=n_0z_0+n_1z_1=|Z|e^{\ii\varphi}$ with $\varphi\in(0,\pi)$,
$\sigma=\sin\varphi$, $c_0=|z_0|\sin(\alpha_0-\varphi)$,
$c_1=|z_1|\sin(\varphi-\alpha_1)$, and
\[
 B_0=D_0^\dagger D_0,\quad L_0=D_0D_0^\dagger,\quad P_1=D_1D_1^\dagger,\quad
 L_1=D_1^\dagger D_1,\quad C=\sum_jA_j^\dagger A_j,\quad G=\sum_jA_jA_j^\dagger .
\]
Here $D_0$ is the differential at the vertex $0$ and $D_1$ the differential
at the vertex $1$; the subscripts are vertices, not degrees. For $r_1=0$
these are complexes as in Theorem~\ref{thm:onevertex}(b) with $i_0=0$, and
for $r_0=0$ as in Theorem~\ref{thm:onevertex}(c) with $i_0=1$.

\begin{proposition}\label{prop:twosidedcoefficient}
The blocks of $\Omega$ for a complex \eqref{eq:twosided} are
\begin{align*}
 \Omega_V&=z_0e^{-B_0}+K_{B_0}(C)-\sum_j\mathcal J_{B_0,L_1}(A_j^\dagger D_1,D_1^\dagger A_j),\\
 \Omega_W&=z_1e^{-P_1}-K_{P_1}(G)+\sum_j\mathcal J_{P_1,L_0}(A_jD_0^\dagger,D_0A_j^\dagger),\\
 \Omega_{U_0}&=\mathcal J_{U_0}-z_0e^{-L_0},\qquad
 \mathcal J_{U_0}:=\sum_j\mathcal J_{L_0,P_1}(D_0A_j^\dagger,A_jD_0^\dagger),\\
 \Omega_{U_1}&=-z_1e^{-L_1}-\mathcal J_{U_1},\qquad
 \mathcal J_{U_1}:=\sum_j\mathcal J_{L_1,B_0}(D_1^\dagger A_j,A_j^\dagger D_1),
\end{align*}
with phases in $(0,\pi)$ on $V,W$, in $(-\pi,0)$ on $U_0$ and in $(\pi,2\pi)$
on $U_1$. If the charges are ordered, then for every metric
\[
 \cD\geq\varphi c_0\tr e^{-L_0}+\sigma\bigl(\varphi\tr\mathcal J_{U_0}-\pi\tr(\mathcal J_{U_0})_-\bigr)
 +(\pi-\varphi)c_1\tr e^{-L_1}+\sigma\bigl((\pi-\varphi)\tr\mathcal J_{U_1}-\pi\tr(\mathcal J_{U_1})_-\bigr).
\]
\end{proposition}

\begin{proof}
Proposition~\ref{prop:coefficient} with $E_0=V\oplus U_0$, $d_0=D_0$,
$E_1=U_1\oplus W$, $d_1=D_1$, $\Delta_0=B_0\oplus L_0$,
$\Delta_1=L_1\oplus P_1$, and the operator
$M_{a_j}=d_1^\dagger\rho(a_j)-\rho(a_j)d_0^\dagger$ of \eqref{eq:DeltaMGamma},
which maps $V\to U_1$ by $D_1^\dagger A_j$ and $U_0\to W$ by
$-A_jD_0^\dagger$. The bound for the block $U_0$ is the argument of
Theorem~\ref{thm:onevertex}(b) with $J$ replaced by the Hermitian matrix
$\mathcal J_{U_0}$. For $U_1$, with $(u_k)$ an eigenbasis of $\Phi_{U_1}$,
eigenvalues $\psi_k\in(\pi,2\pi)$, $q_k=\inner{u_k}{|\Omega_{U_1}|u_k}$,
$\iota_k=\inner{u_k}{\mathcal J_{U_1}u_k}$ and
$e_k=\inner{u_k}{e^{-L_1}u_k}$:
$q_k\sin(\psi_k-\varphi)=\sigma\iota_k+c_1e_k$. If this is nonnegative,
$\psi_k-\varphi\in(\pi-\varphi,\pi]$ and the summand is at least
$(\pi-\varphi)(\sigma\iota_k+c_1e_k)$; otherwise
$\psi_k-\varphi\in(\pi,2\pi-\varphi)$ and it is at least
$(2\pi-\varphi)(\sigma\iota_k+c_1e_k)$. In both cases it is at least
$(\pi-\varphi)(\sigma\iota_k+c_1e_k)-\pi\sigma(\iota_k)_-$.
\end{proof}

In eigenbases of $L_0$, with eigenvalues $\ell_i$, and of $P_1$, with
eigenvalues $p_\mu$ ($p_\mu=0$ on the orthogonal complement of $D_1U_1$),
\[
 (\mathcal J_{U_0})_{ik}=\sum_{j,\mu}(D_0A_j^\dagger)_{i\mu}\overline{(D_0A_j^\dagger)_{k\mu}}\;
 \mathsf f[\ell_i,p_\mu,\ell_k],
\]
with the second divided differences of $e^{-x}$. If all $p_\mu$ vanish this
is the Schur multiplier of \eqref{eq:kernel}; in general the interpolation
point $0$ is replaced by the eigenvalues $p_\mu$.

\begin{theorem}[Ranks at most one]\label{thm:rankone}
Let $\alpha_0\geq\alpha_1$ (ordered charges) and $r_0\leq1$, $r_1\leq1$. Then for every metric
on a complex \eqref{eq:twosided}
\[
 \cD\geq\varphi\bigl(c_0\tr e^{-L_0}+\sigma\tr\mathcal J_{U_0}\bigr)
 +(\pi-\varphi)\bigl(c_1\tr e^{-L_1}+\sigma\tr\mathcal J_{U_1}\bigr)\geq0,
\]
where a term with $r_i=0$ is omitted. The phase potential is nonincreasing
along the phase flow, and strictly decreasing if the charges are strictly
ordered, $n_0,n_1\geq1$ and $r_0+r_1\geq1$.
\end{theorem}

\begin{proof}
For $r_0=1$, $\mathcal J_{U_0}$ is the nonnegative number
\[
 \sum_j\int_{\substack{x,y\geq0\\x+y\leq1}}e^{-(x+y)L_0}\,D_0A_j^\dagger e^{-(1-x-y)P_1}A_jD_0^\dagger\,dx\,dy ;
\]
likewise $\mathcal J_{U_1}\geq0$ for $r_1=1$. Apply
Proposition~\ref{prop:twosidedcoefficient}. If $\alpha_0>\alpha_1$ and
$n_0,n_1\geq1$, then $c_0,c_1>0$.
\end{proof}

\begin{theorem}[Ranks $(2,1)$]\label{thm:twosidedfailure}
Let $z_0=-1+\ii$, $z_1=1+\ii$, and let $V=\CC^3$, $W=\CC^2$, $U_0=\CC^2$,
$U_1=\CC$ with the standard metrics,
\[
 D_0=\begin{pmatrix}0&1&0\\0&0&4\end{pmatrix},\qquad
 D_1=\begin{pmatrix}0\\3\end{pmatrix},\qquad
 A_1=\begin{pmatrix}1&0&0\\0&4&100\end{pmatrix},\qquad A_j=0\ (j\geq2).
\]
The cohomology is the representation $\CC\xrightarrow{1}\CC$, which is
stable, and the induced metric on it is harmonic. For these metrics
\[
 \cD=-1.7280352\ldots<0,\qquad\text{that is,}\qquad
 \frac d{dt}\Ima\bigl(e^{-\ii\varphi}\cS_{\CC}\bigr)=+1.728\ldots>0 .
\]
\end{theorem}

\begin{proof}
The complex is the orthogonal direct sum of the representation
$\CC\xrightarrow1\CC$, whose two blocks $z_0+1$ and $z_1-1$ equal $\ii$ and
have phase $\varphi=\frac\pi2$, and of the acyclic complex
$(0\to U_1)\to(\CC^2\xrightarrow{\rho}\CC)\to(U_0\to0)$ with
$\rho=(\rho_1,\rho_2)=(4,100)$ of Example~\ref{ex:nonnormal}. With
$(\ell_1,\ell_2)=(1,16)$ the eigenvalues of $B_0$ on $\CC^2$ and $p=9$,
Proposition~\ref{prop:twosidedcoefficient} gives for the acyclic complex the
blocks
\begin{align*}
 (\Omega_V)_{ik}&=z_0e^{-\ell_i}\delta_{ik}+\rho_i\rho_k\bigl(-\mathsf f[\ell_i,\ell_k]-p\,\mathsf f[\ell_i,p,\ell_k]\bigr),\\
 \Omega_W&=(z_1-|\rho|^2)e^{-p}+\sum_i\rho_i^2\ell_i\,\mathsf f[p,\ell_i,p],\\
 (-\Omega_{U_0})_{ik}&=z_0e^{-\ell_i}\delta_{ik}-\sqrt{\ell_i\ell_k}\,\rho_i\rho_k\,\mathsf f[\ell_i,p,\ell_k],\\
 -\Omega_{U_1}&=z_1e^{-p}+p\sum_i\rho_i^2\,\mathsf f[p,\ell_i,p],
\end{align*}
with the divided differences of $e^{-x}$. The four block contributions, on
$V$, $W$, $U_0$ and $U_1$, are $3.4635\ldots$, $2.0009\ldots$,
$-10.6257\ldots$ and $3.4332\ldots$; they and their sum were enclosed with
interval arithmetic \cite{joh17},\footnote{The program
\texttt{theorem\_6\_4\_dissipation.py}, which uses the library Arb, and its
output accompany this paper as supplementary files. The program also checks
that $\tr\Omega=0$ for the acyclic complex and compares the blocks, in
floating point, with an evaluation of the definition \eqref{eq:Omegadef},
with the exponential taken in $\mathcal A(K_1)\otimes\End(E)$.} with radius below
$10^{-18}$. The $2\times2$ polar decompositions are explicit, and the phases
on $U_0$ are $-3.1005\ldots$ and $-0.0347\ldots$
\end{proof}

\begin{corollary}\label{cor:classification}
For the charges $(-1+\ii,1+\ii)$ and $n_0=n_1\geq1$, the phase potential is
nonincreasing along the phase flow for all metrics on all complexes
\eqref{eq:twosided} with ranks $(r_0,r_1)$ if and only if $\min(r_0,r_1)=0$
or $r_0=r_1=1$.
\end{corollary}

\begin{proof}
Here $\varphi=\frac\pi2$ and $\alpha_1<\varphi<\alpha_0$. If $r_0=r_1=0$ the
complex is a representation and $\cD\geq0$ by
Proposition~\ref{prop:kempfness}. If
exactly one of $r_0,r_1$ vanishes, this is Theorem~\ref{thm:onevertex}(b) or
(c) for $Q=K_m$; if $r_0=r_1=1$ it is Theorem~\ref{thm:rankone}. Conversely, by Theorem~\ref{thm:duality}, which
exchanges $r_0$ and $r_1$ and fixes these charges, we may assume $r_0\geq2$
and $r_1\geq1$. Take the orthogonal direct sum of $n_0$ copies of
$\CC\xrightarrow1\CC$, of the acyclic complex of
Theorem~\ref{thm:twosidedfailure}, and of contractible summands at the vertex
$0$ in degrees $0,1$ and at the vertex $1$ in degrees $-1,0$, of ranks
$r_0-2$ and $r_1-1$, with differentials $x\,\id$. By
Proposition~\ref{prop:contractible} the latter contribute
$\pi(r_0+r_1-3)e^{-x^2}$ to $\cD$, which tends to $0$ as $x\to\infty$.
\end{proof}

\begin{remark}\label{rem:interpolation}
The example of Theorem~\ref{thm:twosidedfailure} comes from the following
computation. Let $r_0=2$, let $P_1$ have one nonzero eigenvalue $p$, and let
there be one nonzero arrow, which maps into the corresponding line. Then
$(\mathcal J_{U_0})_{ik}=v_i\overline{v_k}\,\mathsf f[\ell_i,p,\ell_k]$ for a
vector $v\in\CC^2$, and for a $2\times2$ matrix of this form the supremum over
$v$ of $\tr(\mathcal J_{U_0})_-/\tr\mathcal J_{U_0}$ is the positive part of
$\frac12\bigl(\mathsf f[\ell_1,p,\ell_2]\,(\mathsf f[\ell_1,p,\ell_1]\mathsf f[\ell_2,p,\ell_2])^{-1/2}-1\bigr)$,
attained for $|v_1|^2\mathsf f[\ell_1,p,\ell_1]=|v_2|^2\mathsf f[\ell_2,p,\ell_2]$.
As $\ell_1\to0$ and $\ell_2\to\infty$ this tends to
\[
 \frac12\Bigl(\frac{R(p)\,e^{p/2}}{\sqrt{(e^{-p}-1+p)/p^2}}-1\Bigr),
\]
which is $(\sqrt2-1)/2$ at $p=0$ and tends to infinity with $p$. So
\eqref{eq:kernel} has no analogue that is uniform in $p$.

\end{remark}

Whether the acyclic parts decay along the phase flow of the complexes
\eqref{eq:twosided}, in the sense of Definition~\ref{def:phase}(v), is open.
